% Источник решений: «Семинар 7.pdf», страницы 2--6.
% Условия: «Неделя 7 Математические ожидания.pdf», conditions/sem_07.pdf.
\section{Математическое ожидание и дисперсия}

\noindent Условия:
\href{run:conditions/sem_07.pdf}{семинар 7 (PDF)}.

\begingroup
\everymath{\textstyle}
\tcbset{myboxstyle/.append style={breakable=false}}
\subsection{Задача 1}

\begin{solution}
    \textbf{Геометрическое распределение.}
    \[
        \Pr\{\xi=k\}=(1-p)^{k-1}p,\qquad k\in\N,\quad 0<p<1.
    \]
    По определению математического ожидания
    \[
        \EE\xi=p\sum_{k=1}^{\infty}k(1-p)^{k-1}.
    \]
    Для $|x|<1$ почленно дифференцируем геометрический ряд:
    \[
        \sum_{k=1}^{\infty}kx^{k-1}
        =\left(\sum_{k=1}^{\infty}x^k\right)'
        =\left(\frac{x}{1-x}\right)'
        =\frac1{(1-x)^2}.
    \]
    Подставляя $x=1-p$, получаем
    \[
        \boxed{\EE\xi=p\cdot\frac1{p^2}=\frac1p.}
    \]

    \textbf{Второй момент и дисперсия.}
    % В рукописи в начальной сумме второго момента пропущен множитель p.
    \[
        \begin{aligned}
            \EE\xi^2
            &=p\sum_{k=1}^{\infty}k^2(1-p)^{k-1}\\
            &=p\sum_{k=2}^{\infty}k(k-1)(1-p)^{k-1}
              +p\sum_{k=1}^{\infty}k(1-p)^{k-1}\\
            &=p(1-p)\sum_{k=2}^{\infty}k(k-1)(1-p)^{k-2}
              +\frac1p.
        \end{aligned}
    \]
    Ещё раз дифференцируем ряд:
    \[
        \sum_{k=2}^{\infty}k(k-1)x^{k-2}
        =\left(\frac1{(1-x)^2}\right)'
        =\frac2{(1-x)^3}.
    \]
    Следовательно,
    \[
        \EE\xi^2=p(1-p)\frac2{p^3}+\frac1p
        =\frac{2-p}{p^2},
    \]
    \[
        \boxed{\DD\xi=\EE\xi^2-(\EE\xi)^2
        =\frac{2-p}{p^2}-\frac1{p^2}=\frac{1-p}{p^2}.}
    \]
    При $p=1$ величина $\xi=1$ почти наверное; те же формулы дают
    $\EE\xi=1$ и $\DD\xi=0$.
\end{solution}

\clearpage
\subsection{Задача 2}

\begin{solution}
    \textbf{Абсолютно непрерывное распределение.}
    \[
        p_\xi(x)=\begin{cases}
            3/x^4,&x\geq1,\\
            0,&x<1.
        \end{cases}
        \qquad \eta=\frac1\xi.
    \]
    \begin{center}
        \begin{tikzpicture}[x=1.3cm,y=0.6cm,>=Stealth]
            \draw[->] (-0.5,0)--(4.4,0) node[right] {$x$};
            \draw[->] (0,-0.25)--(0,3.6) node[above] {$p_\xi(x)$};
            \draw[dashed] (0,3)--(1,3)--(1,0);
            \draw[blue,very thick] (-0.45,0)--(1,0);
            \draw[blue,very thick,domain=1:4.1,samples=100]
                plot (\x,{3/(\x^4)});
            \filldraw[blue,fill=solution_back] (1,0) circle (1.7pt);
            \fill[blue] (1,3) circle (1.7pt);
            \node[below left] at (0,0) {$0$};
            \node[below] at (1,0) {$1$};
            \node[left] at (0,3) {$3$};
        \end{tikzpicture}
    \end{center}

    \textbf{Математическое ожидание и дисперсия $\xi$.}
    \[
        \EE\xi=\int_1^{+\infty}x\frac3{x^4}\,\diff x
        =\left[-\frac3{2x^2}\right]_1^{+\infty}=\frac32.
    \]
    % Исправлены пределы и первообразная во втором моменте.
    \[
        \EE\xi^2=\int_1^{+\infty}x^2\frac3{x^4}\,\diff x
        =\int_1^{+\infty}\frac3{x^2}\,\diff x
        =\left[-\frac3x\right]_1^{+\infty}=3.
    \]
    \[
        \boxed{\EE\xi=\frac32,\qquad
        \DD\xi=3-\frac94=\frac34.}
    \]

    \textbf{Математическое ожидание и дисперсия $\eta=1/\xi$.}
    По формуле математического ожидания функции случайной величины
    \[
        \EE\eta=\EE\frac1\xi
        =\int_1^{+\infty}\frac1x\frac3{x^4}\,\diff x
        =\left[-\frac3{4x^4}\right]_1^{+\infty}=\frac34.
    \]
    % В рукописи вычислено только E eta; далее завершён расчёт D eta
    % по требованию условия тем же способом.
    \[
        \EE\eta^2=\int_1^{+\infty}\frac1{x^2}\frac3{x^4}\,\diff x
        =\left[-\frac3{5x^5}\right]_1^{+\infty}=\frac35.
    \]
    \[
        \boxed{\EE\eta=\frac34,\qquad
        \DD\eta=\frac35-\left(\frac34\right)^2=\frac3{80}.}
    \]
\end{solution}

\clearpage
\subsection{Задача 3}

\begin{solution}
    \textbf{Нормальное распределение.}
    Пусть $\xi\sim N(a;\sigma^2)$, где $a\in\R$, $\sigma>0$:
    \[
        p_\xi(x)=\frac1{\sqrt{2\pi}\,\sigma}
        \exp\left(-\frac{(x-a)^2}{2\sigma^2}\right).
    \]
    Положим $\eta=(\xi-a)/\sigma$. Тогда $\eta\sim N(0;1)$,
    \[
        p_\eta(x)=\frac1{\sqrt{2\pi}}e^{-x^2/2},
        \qquad \xi=\sigma\eta+a.
    \]
    По свойствам линейных преобразований
    \[
        \EE\xi=\sigma\EE\eta+a,
        \qquad \DD\xi=\sigma^2\DD\eta.
    \]

    Подынтегральная функция в следующем абсолютно сходящемся интеграле
    нечётна, поэтому
    \[
        \EE\eta=\frac1{\sqrt{2\pi}}
        \int_{-\infty}^{+\infty}xe^{-x^2/2}\,\diff x=0.
    \]
    Следовательно, $\DD\eta=\EE\eta^2$.
    Вычислим второй момент интегрированием по частям:
    \[
        u(x)=x,\quad u'(x)=1,
        \qquad v'(x)=xe^{-x^2/2},\quad v(x)=-e^{-x^2/2}.
    \]
    \[
        \begin{aligned}
            \EE\eta^2
            &=\frac1{\sqrt{2\pi}}
              \int_{-\infty}^{+\infty}x^2e^{-x^2/2}\,\diff x\\
            &=\frac1{\sqrt{2\pi}}
              \left(\left[-xe^{-x^2/2}\right]_{-\infty}^{+\infty}
              +\int_{-\infty}^{+\infty}e^{-x^2/2}\,\diff x\right)
              =1.
        \end{aligned}
    \]
    Граничный член равен нулю, а оставшийся интеграл с множителем
    $1/\sqrt{2\pi}$ равен $1$ по нормировке стандартной нормальной плотности.
    Таким образом,
    \[
        \boxed{\EE\xi=a,\qquad \DD\xi=\sigma^2.}
    \]
\end{solution}
\endgroup
